\documentclass[a4paper,12pt]{article}

% Eingabekodierung: UTF-8 (ermöglicht direkte Eingabe von Umlauten)
\usepackage[utf8]{inputenc}
% Schriftkodierung: T1 für bessere Trennung und Akzentdarstellung
\usepackage[T1]{fontenc}
% Deutsche Sprachunterstützung (Trennungen, typografische Regeln, Datum etc.)
\usepackage[ngerman]{babel}
% Seitenlayout / Ränder
\usepackage{geometry}

\usepackage{xcolor}    % nur für Blau (optional)




% AMS-Erweiterungen für Mathematik: Umgebungen, Symbole und Schriften
\usepackage{amsmath, amssymb, amsfonts}
% Zusätzliche Mathe-Tools (Erweiterungen und Komfortfunktionen zu amsmath)
\usepackage{mathtools}

\usepackage{trsym}   % definiert \TransformVert, \InversTransformVert, \TransformHoriz, ...

\usepackage{tikz} % TikZ für Zeichnungen (tikzpicture, \draw, \fill, ...)




\geometry{margin=2.5cm}

\title{Darstellung von Formeln und Gleichungssystemen in \LaTeX}
\author{}
\date{}

\begin{document}

\maketitle

\tableofcontents
\newpage

% ==============================================================
\section{Formeln im Fließtext und in eigener Zeile}

Formeln können im Fließtext oder abgesetzt dargestellt werden.

\subsection{Formel im Fließtext}

Beispiel einer Formel im Text:

\noindent Die Energie eines ruhenden Körpers wird durch $E = mc^2$ beschrieben...

\subsection{Formel in eigener Zeile (Display-Math)}

Die gleiche Formel abgesetzt und zentriert:

\[
E = mc^2
\]

\noindent oder mit automatischer Nummerierung:

\begin{equation}
E = mc^2
\end{equation}

\noindent Man kann mit einem Label darauf verweisen, zum Beispiel Gleichung~\eqref{eq:energie}:

\begin{equation}
E = mc^2
\label{eq:energie}
\end{equation}

% ==============================================================
\section{Gleichungssysteme}

Gleichungssysteme können auf verschiedene Arten dargestellt werden:
entweder als Block mit geschweiften Klammern oder als ausgerichtete Zeilen.

\subsection{Einfaches Gleichungssystem (ohne Ausrichtung)}

\[
\begin{cases}
x + y = 5 \\
x - y = 1
\end{cases}
\]

\subsection{Ausgerichtetes Gleichungssystem (mit \texttt{align})}

Hier werden die Gleichheitszeichen untereinander ausgerichtet:

\begin{align}
x + y &= 5 \\
x - y &= 1
\end{align}

\subsection{Gleichungssystem mit Beschriftung (Label)}

\noindent Das Label ermöglicht späteres Referenzieren.  
\noindent Beispiel: Siehe Gleichungssystem~\eqref{eq:gleichungssystem}.

\begin{align}
a + b &= 3 \\
2a - b &= 4
\label{eq:gleichungssystem}
\end{align}

\subsection{Ausrichtung mit mehreren Spalten (z.\,B. Zwischenschritte)}

\begin{align*}
x + y &= 5  && \text{(erste Gleichung)}\\
x - y &= 1  && \text{(zweite Gleichung)}\\
\Rightarrow\; 2x &= 6 && \text{(Addition)}\\
\Rightarrow\; x &= 3 && \text{(Division durch 2)}
\end{align*}

\subsection{Gleichungssystem mit einer gemeinsamen Nummer (split)}

\begin{equation}
\label{eq:split-system}
\begin{split}
x + y &= 5 \\
x - y &= 1
\end{split}
\end{equation}

\noindent Das obige System~\eqref{eq:split-system} wird als eine Einheit nummeriert.

% ==============================================================
\section{Laplace-Transformation}
\[
\mathcal{L}\{f(t)\} = F(s)
\]

\[
\begin{aligned}
f(t)\\
\TransformVert\\
F(s)
\end{aligned}
\]

% ==============================================================
\section{Matrix}
\noindent Ein Beispiel aus der Statistik:

% Darstellung mit runder Klammern
\[
A = \begin{pmatrix}
  a_{11} & a_{12} & a_{13} \\
  a_{21} & a_{22} & a_{23} \\
  a_{31} & a_{32} & a_{33}
\end{pmatrix}
\]

% Darstellung mit eckigen Klammern
\[
B = \begin{bmatrix}
  1 & 2 & 3 \\
  4 & 5 & 6 \\
  7 & 8 & 9
\end{bmatrix}
\]

% Darstellung mit geschwungener Klammern
\[
C = \left\{
\begin{matrix}
c_{11} & c_{12} & c_{13} \\
c_{21} & c_{22} & c_{23} \\
c_{31} & c_{32} & c_{33}
\end{matrix}
\right\}
\]

% ==============================================================
\section{Sonderzeichen und Operatoren}

\[
\begin{aligned}
x^2, \quad a_i, \quad \frac{a}{b}, \quad \sqrt{x}, \quad \sqrt[3]{y}, \\ % Potenz, Index, Bruch, Quadratwurzel, Kubikwurzel
\sum_{i=1}^n i^2, \quad \int_0^\infty e^{-x}\,dx, \quad % Summe, uneigentliches Integral (mit kleinem Abstand vor dx)
\lim_{x\to0} \frac{\sin x}{x}, \\ % Limes
\vec{a}, \quad \overrightarrow{AB}, \quad |x|, \quad \lvert x \rvert % Vektornotationen, Betrag (zwei Schreibweisen)
\end{aligned}
\]

% ==============================================================
\section{Kleine griechische Buchstaben im Fließtext}

Hier einige Beispiele für kleine griechische Buchstaben im Text:

\noindent $\alpha, \beta, \gamma, \delta, \epsilon, \zeta, \eta, \theta, \iota, \kappa, \lambda, \mu, \nu, \xi, \pi, \rho, \sigma, \tau, \phi, \chi, \psi, \omega$.

\noindent Beispielhafte Verwendung im Kontext:

\noindent Die Dichte $\rho$ eines Körpers und der Winkel $\theta$ sind wichtige Größen in der Physik.

% ==============================================================
\section{Große griechische Buchstaben in Gleichungen}

In der Mathematik werden auch große griechische Buchstaben verwendet:

\[
\Gamma + \Delta + \Theta = \Lambda + \Pi + \Sigma + \Omega
\] 



% ==============================================================
\section{Beispielhafte Formeln mit griechischen Symbolen}


\subsection{Beispiel aus der Trigonometrie}

\[
\sin(\alpha) + \cos(\beta) = 1
\]

\subsection{Beispiel aus der Physik}

\[
F = m \cdot g = m \cdot 9.81~\frac{\text{m}}{\text{s}^2}
\]

\noindent mit dem Winkelbezug $\theta$:
\[
F_\text{Hang} = F \cdot \sin(\theta)
\]

% ==============================================================
\section{Übersicht: Griechische Buchstaben in \LaTeX}

\begin{center}
\begin{tabular}{ll|ll}
\textbf{Klein} & \textbf{Code} & \textbf{Groß} & \textbf{Code} \\
\hline
$\alpha$ & \verb|\alpha| & $\Gamma$ & \verb|\Gamma| \\
$\beta$ & \verb|\beta| & $\Delta$ & \verb|\Delta| \\
$\gamma$ & \verb|\gamma| & $\Theta$ & \verb|\Theta| \\
$\delta$ & \verb|\delta| & $\Lambda$ & \verb|\Lambda| \\
$\epsilon$ & \verb|\epsilon| & $\Xi$ & \verb|\Xi| \\
$\zeta$ & \verb|\zeta| & $\Pi$ & \verb|\Pi| \\
$\eta$ & \verb|\eta| & $\Sigma$ & \verb|\Sigma| \\
$\theta$ & \verb|\theta| & $\Upsilon$ & \verb|\Upsilon| \\
$\lambda$ & \verb|\lambda| & $\Phi$ & \verb|\Phi| \\
$\mu$ & \verb|\mu| & $\Psi$ & \verb|\Psi| \\
$\nu$ & \verb|\nu| & $\Omega$ & \verb|\Omega| \\
$\xi$ & \verb|\xi| & & \\
$\pi$ & \verb|\pi| & & \\
$\rho$ & \verb|\rho| & & \\
$\sigma$ & \verb|\sigma| & & \\
$\tau$ & \verb|\tau| & & \\
$\phi$ & \verb|\phi| & & \\
$\chi$ & \verb|\chi| & & \\
$\psi$ & \verb|\psi| & & \\
$\omega$ & \verb|\omega| & & \\
\end{tabular}
\end{center}

% ==============================================================
\section{Mengenzeichen}

Die wichtigsten Zahlenmengen der Mathematik werden mit Doppelstrich-Buchstaben dargestellt.  
Dazu verwendet man in \LaTeX{} den Befehl \verb|\mathbb|{} (bereitgestellt durch das Paket \texttt{amssymb}).

\[
\mathbb{N}, \mathbb{Z}, \mathbb{Q}, \mathbb{R}, \mathbb{C}
\]

\begin{itemize}
    \item $\mathbb{N}$ – Menge der natürlichen Zahlen
    \item $\mathbb{Z}$ – Menge der ganzen Zahlen
    \item $\mathbb{Q}$ – Menge der rationalen Zahlen
    \item $\mathbb{R}$ – Menge der reellen Zahlen
    \item $\mathbb{C}$ – Menge der komplexen Zahlen
\end{itemize}
\end{document}
